\fichatitulo{25 · REVISÃO}{Preprint v1 · 2024}{Panorama de grandes modelos de linguagem}
{\small Minaee et al. \textit{Large Language Models: A Survey}. Preprint v1 · 2024. \href{https://arxiv.org/pdf/2402.06196v1}{Fonte consultada}.\par}

\campo{Problema e objetivo}
Como organizar famílias, desenvolvimento e avaliação dos LLMs?

\campo{Principal contribuição · síntese interpretativa}
Reúne características e limitações de famílias de modelos e organiza recursos para sua construção, utilização e ampliação.

\campo{Método / natureza da evidência}
Revisão de literatura técnica; apresenta modelos, dados, avaliação e problemas em aberto.

\campo{Resultado ou argumento principal}
Oferece um mapa conceitual para situar abordagens e vocabulário. Não constitui um novo experimento comparativo de RAG.

\campo{Excertos-chave · seleção literal}
\begin{itemize}
\excerto{1}{characteristics, contributions and limitations}{Resumo.}
\excerto{2}{LLM evaluation metrics}{Resumo.}
\end{itemize}

\campo{Limitações e análise crítica}
Análise da ficha: a cobertura reflete a versão de fevereiro de 2024 e não equivale a um inventário atualizado. Resultados de modelos específicos devem ser citados a partir dos estudos primários.

\campo{Aproveitamento na dissertação · proposta}
Fundamentação: selecionar conceitos necessários e encaminhar afirmações empíricas às fontes originais.

\registro{Fonte consultada nesta preparação: Resumo, mapa de organização e trechos de desafios; consulta parcial. Chave: \texttt{minaee2024llmSurvey}.}
